% Descrever brevemente o produto e incluir os benefícios esperados, seu propósito principal e quais necessidades ele busca atender.

% Página 67 da ISO 29148

O sistema proposto, denominado \textbf{SpinFlow}, consiste em uma plataforma digital destinada ao apoio de academias na organização e gerenciamento de aulas de \textit{spinning}.

O principal propósito do \textbf{SpinFlow} é melhorar o processo de agendamento e controle das aulas, oferecendo suporte ao gerenciamento de alunos, salas e recursos associados, como bicicletas e estrutura física. Além disso, o sistema contempla funcionalidades auxiliares relacionadas à personalização das aulas, como categorias musicais e artistas, bem como apoio ao controle de manutenção dos equipamentos.

A solução busca atender necessidades recorrentes identificadas no contexto das academias, tais como dificuldades no controle de vagas, organização das turmas, gerenciamento de recursos físicos e ausência de ferramentas estruturadas em ambientes que ainda utilizam processos manuais ou sistemas limitados.

Como benefícios esperados, destacam-se a melhoria na organização das aulas, maior controle sobre a capacidade das salas, otimização do uso dos recursos disponíveis, apoio à gestão de manutenção e aprimoramento da experiência dos alunos.

O escopo do sistema abrange o gerenciamento das seguintes entidades principais: alunos, salas, categorias de música, artistas/bandas, tipos de manutenção e fabricantes, permitindo uma abordagem integrada entre aspectos operacionais e administrativos.